\id{МРНТИ 65.43.29}{}

\begin{header}
\swa{Д.Н.~Бакытжан, Э.Б.~Аскарбеков, И.В.~Новикова, А.К.~Кекибаева, Г.И.~Байгазиева}{ОТАНДЫҚ СЕЛЕКЦИЯЛЫ БИДАЙ ДӘНІНЕН СЫРАЛЫҚ УЫТ ӨНДІРУГЕ АРНАЛҒАН УЫТ ӨСІРУ РЕЖИМДЕРІН ТАҢДАУ}

\tsp{1}Д.Н.~Бакытжан\envelope,
\tsp{2}Э.Б.~Аскарбеков,
\tsp{3}И.В.~Новикова,
\tsp{1}А.К.~Кекибаева,
\tsp{1}Г.И.~Байгазиева
\end{header}

\begin{affil}
\tsp{1}Алматы технологиялық университеті, Алматы, Қазақстан,

\tsp{2}Қ.Құлажанов атындағы Қазақ технология және бизнес университеті, Астана, Қазақстан,

\tsp{3}Воронеж мемлекеттік инженерлік технологиялар университеті, Воронеж, Ресей

\corrauthor{Корреспондент-автор: didar\_99\_kz@mail.ru}
\end{affil}

Мақалада Қазақстан Республикасының отандық селекциясымен шығарылған
бидай сорттарынан бидай уытын өндірудің теориялық және эксперименттік
аспектілері қарастырылып, оны сыра қайнату саласында қолданудың мақсатқа
сәйкестігі негізделді. Зерттеудің өзектілігі жергілікті дәнді дақыл
шикізатын пайдалану арқылы сыра қайнатудың шикізат базасын кеңейту және
импорттық уытқа тәуелділікті төмендету қажеттілігімен айқындалады.
Зерттеу нысандары ретінде дәннің физика-химиялық және технологиялық
сипаттамалары бойынша ерекшеленетін қазақстандық селекциялы жұмсақ бидай
сорттары - Матай, Ертіс және Тәуелсіздік таңдап алынды. Жұмыста
ғылыми-техникалық және нормативтік әдебиеттерді талдау әдістері, сонымен
қатар дән мен уыттың сапасын бағалаудың жалпы қабылданған
физика-химиялық, биохимиялық және технологиялық әдістері қолданылды.
Эксперименттік бөлім бидайды уыт өсірудің әртүрлі
температуралық-уақыттық режимдерін таңдау мен салыстыруды, ылғал
жинақталуының динамикасын, азотты зат\-тар мөлшерінің және бос аминді
азоттың өзгерісін бақылауды қамтыды. Зерттелген бидай сорттарының
барлығы қанағаттанарлық физиологиялық белсенділікке ие екені анықталды
және оларды уыт өндірісі үшін шикізат ретінде қарастыруға болады. Бидай
уытын алу үшін ең қолайлы Тәуелсіздік сорты болып шықты, ол ақуыз
мөлшерінің және дәннің шынылығының орташа деңгейімен сипатталады. Уыт
өсіру температурасы ақуыздық гидролиздің қарқындылығына және ашытқылар
үшін уыттың қоректік құндылығын айқындайтын еркін аминді азоттың
жиналуына елеулі әсер ететіні көрсетілді. Алынған нәтижелер
технологиялық көрсеткіштері сыра қайнату саласының талаптарына сәйкес
келетін бидай уытын өндіру мүмкіндігін растайды және отандық шикізат
негізінде бидай сырасын өндіру технологияларын әзірлеуде пайдаланылуы
мүмкін.

{\bfseries Түйін сөздер:} сыра қайнату, сыралық уыт, технология, бидай, уыт
өсіру, аминді азот.

\begin{header}
ПОДБОР РЕЖИМОВ СОЛОДОРАЩЕНИЯ ДЛЯ ПРОИЗВОДСТВА ПИВОВАРЕННОГО СОЛОДА ИЗ ЗЕРНА ПШЕНИЦЫ ОТЕЧЕСТВЕННОЙ СЕЛЕКЦИИ

\tsp{1}Д.~Бакытжан\envelope,
\tsp{2}Э.Б.~Аскарбеков,
\tsp{3}И.В.~Новикова,
\tsp{1}А.К.~Кекибаева,
\tsp{1}Г.И.~Байгазиева
\end{header}

\begin{affil}
\tsp{1}Алматинский технологический университет, Алматы, Казахстан,

\tsp{2}Казахский университет технологии и бизнеса им.Кулажанова К.С., Астана, Казахстан,

\tsp{3}Воронежский государственный университет инженерных технологий, Воронеж, Россия,

е-mail: didar\_99\_kz@mail.ru
\end{affil}

В статье рассмотрены теоретические и экспериментальные аспекты
производства пшеничного солода из сортов пшеницы отечественной селекции
Республики Казахстан и обоснована целесообразность его применения в
пивоваренной отрасли. Актуальность исследования обусловлена
необходимостью расширения сырьевой базы пивоварения за счёт
использования местного зернового сырья и снижения зависимости от
импортного солода. В качестве объектов исследования выбраны сорта мягкой
пшеницы казахстанской селекции: Матай, Ертіс и Тәуелсіздік,
различающиеся по физико-химическим и технологическим характеристикам
зерна. В работе применены методы анализа научно-технической и
нормативной литературы, а также общепринятые физико-химические,
биохимические и технологические методы оценки качества зерна и солода.
Экспериментальная часть включала подбор и сравнение различных
температурно-временных режимов солодоращения пшеницы с контролем
динамики влагонакопления, изменения содержания азотистых веществ и
свободного аминного азота. Установлено, что все исследуемые сорта
пшеницы обладают удовлетворительной физиологической активностью и могут
рассматриваться в качестве сырья для солодовенного производства.
Наиболее перспективным для получения пшеничного солода оказался сорт
Тәуелсіздік, характеризующийся умеренным содержанием белка и
стекловидности. Показано, что температура солодоращения оказывает
существенное влияние на интенсивность белкового гидролиза и накопление
свободного аминного азота, определяющего питательную ценность солода для
дрожжей. Полученные результаты подтверждают возможность получения
пшеничного солода с технологическими показателями, соответствующими
требованиям пивоваренной отрасли, и могут быть использованы при
разработке технологий производства пшеничного пива на основе
отечественного сырья.

{\bfseries Ключевые слова:} пивоварение, солод, технология, пшеница,
солодоращение, аминный азот.

\begin{header}
SELECTION OF MALTING REGIMES FOR THE PRODUCTION OF BREWING MALT FROM WHEAT GRAIN OF DOMESTIC BREEDING

\tsp{1}D.~Bakytzhan\envelope,
\tsp{2}E.B.~Askarbekov,
\tsp{3}I.V.~Novikova,
\tsp{1}A.K.~Kekibayeva,
\tsp{1}G.I.~Baigazieva
\end{header}

\begin{affil}
\tsp{1}Almaty Technological University, Almaty, Kazakhstan,

\tsp{2}K.S. Kulazhanov Kazakh University of Technology and Business, Astana Kazakhstan,

\tsp{3}Voronezh State University of Engineering Technologies, Voronezh, Russia,

e-mail: didar\_99\_kz@mail.ru
\end{affil}

The article examines the theoretical and experimental aspects of wheat
malt production from wheat varieties of domestic breeding of the
Republic of Kazakhstan and substantiates the feasibility of its
application in the brewing industry. The relevance of the study is
обусловлена the need to expand the raw material base of brewing through
the use of local grain resources and to reduce dependence on imported
malt. The objects of the study were soft wheat varieties of Kazakh
Breeding-Matai, Ertis, and Tauelsizdik-which differ in their
physicochemical and technological grain characteristics. The research
employed methods of analysis of scientific, technical, and regulatory
literature, as well as generally accepted physicochemical, biochemical,
and technological methods for assessing the quality of grain and malt.
The experimental part included the selection and comparison of various
temperature-time malting regimes for wheat, with monitoring of moisture
uptake dynamics, changes in nitrogenous substances, and free amino
nitrogen content. It was established that all the studied wheat
varieties exhibit satisfactory physiological activity and can be
considered as raw materials for malting production. The Tauelsizdik
variety proved to be the most promising for wheat malt production, being
characterized by moderate protein content and vitreousness. It was shown
that malting temperature has a significant effect on the intensity of
protein hydrolysis and the accumulation of free amino nitrogen, which
determines the nutri\-tional value of malt for yeast. The obtained results
confirm the possibility of producing wheat malt with techno\-logical
indicators meeting the requirements of the brewing industry and can be
used in the development of technologies for wheat beer production based
on domestic raw materials.

{\bfseries Keywords:} brewing, malt, technology, wheat, malting, free amino
nitrogen (FAN).

\begin{multicols}{2}
{\bfseries Кіріспе.} Сыра қайнату - тағам өнеркәсібінің ең көне, әрі
технологиялық тұрғыдан күрделі салаларының бірі. Бұл саланың негізінде
ашытылатын қанттардың, ферменттердің және дәмдік-хош иістік
қосылыстардың негізгі көзі ретінде уытты пайдалану жатыр {[}1,2{]}. Арпа
уытымен қатар, бидай уыты да сыраның кең ассортиментін өндіруде, әсіресе
өзіндік органолептикалық сипаттамаларымен ерекшеленетін жоғары
ашытылатын сорттарды дайындауда маңызды рөл атқарады {[}3{]}.

Бидай уыты бидай дәнін бақыланатын жағдайда өсіру және кейіннен кептіру
арқылы алынады, соның нәтижесінде ферменттік жүйелер, ең алдымен
амилазалар мен протеазалар белсенденеді {[}4{]}. Бұл ферменттер уытты
ысқылау процесінде дәннің крахмалы мен ақуызды заттарының гидролизін
қамтамасыз етіп, суслоның құрамын қалыптастырады және дайын сыраның
сапасын айқындайды. Арпадан айырмашылығы, бидай қауыссыз дән, бұл
өңдеудің технологиялық параметрлеріне, затордың сүзілгіштік қасиеттеріне
және сусынның құрылымына әсер етеді {[}5{]}.

Сыра қайнатуда бидай уытын қолдану жұмсақ дәмнің қалыптасуына, көбіктің
жоғары түзілу қабілетіне және бидай сырасына тән ерекше лайлылықтың
пайда болуына ықпал етеді {[}6{]}. Бидайдағы ақуыздар мен ерігіш азотты
заттардың мөлшерінің жоғары болуы көбіктің тұрақтылығына оң әсер етеді,
алайда уыттандыру және уытты езгілеу режимдерін неғұрлым дәл бақылауды
талап етеді {[}7{]}.

Қазақстан айтарлықтай аграрлық әлеуетке ие мемлекет және Еуразия
өңіріндегі ірі бидай өндірушілердің бірі болып табылады {[}8{]}. Елдің
климаттық жағдайлары мен топырақ-климаттық аймақтарының алуан түрлілігі
крахмал мөлшері жоғары, шынылығы оңтайлы және ақуыз көрсеткіштері
қанағаттанарлық бидай сорттарын өсіруге мүмкіндік береді, бұл оларды уыт
және сыра өндірісі үшін перспективалы шикізатқа айналдырады. Алайда
дәстүрлі түрде Қазақстандық селекциялы бидай негізінен нан пісіру және
экспорт салаларына бағдарланған, ал оны сыра қайнатуда қолдану мәселесі
әлі де жеткілікті деңгейде зерттелмеген {[}9{]}.

Отандық бидай сорттарынан бидай уытын өндірудің экономикалық және
технологиялық маңызы зор. Жергілікті шикізатты пайдалану сыра қайнату
өнеркәсібінің импорттық уытқа тәуелділігін төмендетуге, тасымалдау
шығындарын қысқартуға және жеткізілімдердің тұрақтылығын қамтамасыз
етуге мүмкіндік береді. Сонымен қатар, уыттандыру режимдерін
қазақстандық бидай сорттарының биохимиялық ерекшеліктеріне бейімдеу
алдын ала белгіленген технологиялық қасиеттері бар уыт әзірлеуге жол
ашады.

Ғылыми тұрғыдан алғанда, қазақстандық бидай сорттарының уыттық
қасиеттерін зерттеу сыра қайнатудың шикізат базасын кеңейту және дәнді
өңдеу технологияларын жетілдіру үшін қызығушылық тудырады.
Ферментативтік белсенділікті, экстрактивтілікті, ақуыздардың ерігіштігін
және бидай уыты сапасының өзге де көрсеткіштерін зерделеу оны сыраның
әртүрлі рецептураларында қолданудың мақсатқа сәйкестігін негіздеуге
мүмкіндік береді {[}3{]}.

Осыған байланысты, Қазақстанда өсірілетін бидай сорттарынан бидай уытын
өндіруді дамыту - сыра қайнату саласының инновациялық дамуына ықпал
ететін, аграрлық ресурстарды ұтымды пайдалануды қамтамасыз ететін және
отандық өнімнің ішкі әрі сыртқы нарықтардағы бәсекеге қабілеттілігін
арттыратын перспективалы бағыт болып табылады {[}10{]}.

Ұсынылып отырған зерттеудің мақсаты - отандық селекциялы бидай
сорттарынан сыра қайнатуға қолайлы бидай уытын өндірудің сәйкестігін
негіздеу.

Қойылған мақсатқа қол жеткізу үшін зерттеуде мынадай міндеттерді шешу
көзделеді: уыт өндірісі үшін маңызды қазақстандық селекциялы бидай
дәнінің биохимиялық және технологиялық ерекшеліктерін зерттеу; уыт
өндіруде теориялық негіздерін және процестердің бидай уытының
ферментативтік белсенділігіне әсерін талдау.

Зерттеу гипотезасы Қазақстанда өсірілетін бидай сорттары уыт өсіру
режимдерін оңтайландыру жағдайында сыра қайнату өнеркәсібінің
талаптарына сай жеткілікті ферментативтік белсенділігі бар бидай уытын
алуды қамтамасыз ете алады және бидай сырасының түрлерін өндіруде тиімді
пайдаланылуы мүмкін деген болжамға негізделеді.

{\bfseries Материалдар мен әдістер.} Зерттеу нысандары ретінде отандық
селекциялы бидай сорттары алынды: №1 үлгі -Матай, №2 үлгі - Ертіс, №3
үлгі - Тәуелсіздік, сондай-ақ сыра қайнатуға арналған бидай уытын өндіру
үдерісінің өзі қарастырылды. Бидай сорттары «Қазақ егіншілік және
өсімдік шаруашылығы ғылыми-зерттеу институты» ЖШС-нен (Алматы облысы,
Қазақстан) алынды. Ғылыми зерттеулер Алматы технологиялық
университетінің «Тағам қауіпсіздігі» ғылыми-зерттеу институтында,
сонымен қатар ФГБОУ ВО ВГУИТ (РФ, Воронеж қаласы) сынақ орталығында
жүргізілді.

Жұмыста келесі зерттеу әдістері қолданылды: уыт өсіру және сыра қайнату
бойынша ғылыми-техникалық әрі нормативтік әдебиеттерді талдау және
жалпылау; дән мен уыт сапасының физика-химиялық көрсеткіштерін,
сондай-ақ дәннің уыттық көрсеткіштерін анықтаудың жалпы қабылданған
әдістері.

Бидай дәнінің өнгіштік қабілетін анықтау МЕМСТ 10968-88 стандарты
бойынша жүргізілді. Әдіс қатаң белгіленген температуралық жағдайларда
ылғалды ортада дән өлшендісін өндіруге негізделеді, кейін қалыпты өскін
түзген дәндердің саны есептеледі. Ол үшін орташа сынамадан зақымданған
және толық дамымаған дәндерсіз таза сынама бөлініп алынады. Дәндер сүзгі
қағазына немесе Петри табақшаларына орналастырылады. Қағаз дистилденген
сумен біркелкі ылғалдандырылады. Сынамалар 20 ± 2°C температурада
термостатқа қойылып, өндіру ұзақтығы 5-7 тәулік жүргізіледі. Қалыпты
өскін берген дәндердің саны есептеліп, көрсеткіш формула бойынша
есептеледі.

Дәннің тіршілікке қабілеттілігін анықтау МЕМСТ 10968-88 стандарты
бойынша жүргізілді. Әдіс физиологиялық тыныштық күйінде болатын дәндерді
қоса алғанда, зат алмасу үдерістеріне қабілетті тірі дәндерді анықтауға
негізделген. Дәндер 18--22°C температурадағы суға 24 сағат жібітіледі.
Жібіткеннен кейін жеделдетілген өндіру жүргізіледі немесе ұрықты бояу
әдісі қолданылады. Ұрықтың тіршілік белсенділігі белгілері бар тірі
дәндердің саны есептеледі.

{\bfseries Нәтижелер және талқылау.} «Матай» бидай сорты Қазақстанның
оңтүстік-шығыс өңірінің табиғи-климаттық жағдайларын ескере отырып
шығарылған жұмсақ бидайдың қазақстандық сорттарына жатады. Сорт атауы
Матай өңірімен байланысты, бұл оның ылғалдануы тұрақсыз тау бөктері және
дала аймақтарына бейімделгенін көрсетеді.

Сорт орташа мерзімде пісіп-жетілуімен және өңірдің топырақ-климаттық
жағдайларына жақсы бейімделгіштігімен сипатталады. Дәні жеткілікті
дәрежеде біркелкі қалыптасады, шынылығы қанағаттанарлық және ұнтақтауға
қолайлы қасиеттерімен ерекшеленеді.

«Матай» ең алдымен нан пісіру мақсатында пайдалануға бағытталған және
тұрақты технологиялық сипаттамалары бар ұн алуды қамтамасыз етеді. Сорт
өңірлік егіншілікте кеңінен қолданылады және өнімнің тұрақтылығымен
бағаланады. Бұл сортты сыра қайнату саласында қолдану ерекше қызығушылық
тудырады {[}11{]}.

«Ертіс» сорты Қазақстанның солтүстік және солтүстік-шығыс өңірлеріне
арналған жұмсақ бидай сорттарына жатады. Атауы Ертіс өзенімен
байланысты, бұл оның климаты салқындау аймақтарда аудандастырылғанын
көрсетеді.

Сорт интенсивті даму типімен және тығыз, толық толысқан дән түзуімен
ерекшеленеді. Ақуыз мөлшерінің жоғары болуы және желімшенің сапасының
жақсы болуына байланысты «Ертіс» дәні күшті бидайлар санатына жатады.
Бұл ерекшелік одан сыра қайнатуға арналған уыт өндіру тұрғысынан да
қызығушылық тудырады {[}12{]}.

«Тәуелсіздік» сорты - ұлттық селекциялық бағдарламалардың дамуы
кезеңінде шығарылған Қазақстандық селекцияның салыстырмалы түрде
заманауи сорты. Атауы («Тәуелсіздік») оның тұрақтылық пен әмбебаптыққа
бағдарланғанын білдіреді.

«Тәуелсіздік» сортының дәні орташа ірілігімен, шынылығы мен
ақуыздылығының орташа көрсеткіштерімен сипатталады. Бұл сортты қолданылу
бағыты тұрғысынан әмбебап етеді және оның уыттық қасиеттерін зерттеу
барысында ерекше қызығушылық туғызады {[}13{]}.

Зерттеудің бірінші кезеңінде зерттелетін сорттардың физика-химиялық
сипаттамалары зерделеніп, олардың сыра қайнатуға жарамдылығы тұрғысынан
уыттық қасиеттеріне баға берілді (1-кесте).
\end{multicols}

\tcap{1-кесте. Зерттелетін бидай үлгілерінің сапа көрсеткіштері}
\begin{longtblr}[
  label = none,
  entry = none,
]{
  width = \linewidth,
  colspec = {Q[300]Q[200]Q[56]Q[50]Q[81]},
  cells = {c},
  cells = {font = \small},
  row{1} = {font=\bfseries\small},
  hlines,
  vlines,
}
Көрсеткіш & Сыра қайнатуға арналған арпа (МЕМСТ 5060-2021) & Матай & Ертіс & Тәуелсіздік\\
Ылғалдылық, \% & 14,5 & 13,5 & 13,0 & 13,5\\
Дән натурасы, г/л & 650 және одан жоғары & 760 & 780 & 750\\
1000 дән массасы, г & ≥ 45 & 38 & 42 & 36\\
Шынылығы, \% & 50-60 & 65 & 70 & 55\\
Ақуыз мөлшері, \% & 12 & 12,5 & 14,0 & 11,5\\
Өнгіштік қабілеті, \%, кемінде (жинағаннан кейін 45 күннен ерте емес жеткізілетін дән үшін) & 90-95 & 95 & 95 & 95\\
Тіршілікке қабілеттілігі, \%, кемінде (жинағаннан кейін 45 күнге дейінгі мерзімде жеткізілетін дән үшін) & 95 & 90 & 90 & 90
\end{longtblr}

\begin{multicols}{2}
Дәннің уыттық қасиеттері дәнді суландыру, өндіру және эндоспермнің
ферментативтік модификациясы процестеріне әсер ететін физика-химиялық
және физиологиялық көрсеткіштердің жиынтығымен айқындалады. Салыстыру
үшін эталондық шикізат ретінде сыра қайнатуға арналған арпа
пайдаланылды, оған қойылатын талаптар МЕМСТ 5060-2021 стандартымен
реттеледі.

1-кестеде келтірілген деректерге сәйкес зерттелген бидай сорттарының
уыттық қасиеттеріне салыстырмалы талдау жүргізілді. Бидай дәнінің
ылғалдылығы 13,0-13,5\% аралығында өзгерді, бұл сыра қайнатуға арналған
арпа үшін рұқсат етілетін шекті мәннен (14,5\%) төмен. Мұндай ылғалдылық
деңгейі сақтау кезінде дәннің технологиялық тұрақтылығын қамтамасыз
етіп, микробиологиялық үдерістердің даму қаупін төмендетеді, бұл
шикізатты уыттандыруға дайындау барысында оң фактор болып табылады.

Матай, Ертіс және Тәуелсіздік сорттарының дән натурасы 750-780 г/л
болды, бұл сыра қайнатуға арналған арпаға қойылатын ең төменгі
талаптардан (кемінде 650 г/л) едәуір жоғары. Аталған көрсеткіштің ең
жоғары мәні Ертіс сортында (780 г/л) байқалды, бұл эндоспермнің жақсы
толысқанын және тығыздығы жоғары екенін көрсетеді. Дән натурасының
жоғары болуы, бір жағынан, экстракттың әлеуетті шығымын арттыруға ықпал
етсе, екінші жағынан, дәннің қаттылығының жоғары екенін білдіріп, уыт
өсіру кезінде эндоспермнің модификациялану үдерістерін күрделендіруі
мүмкін.

Зерттелген бидай сорттарының 1000 дән массасы 36-42 г аралығында болды,
бұл сыра қайнатуға арналған арпа үшін оңтайлы мәннен (≥ 45 г) төмен.
Дәннің салыстырмалы түрде ірірек болуы Ертіс сортына тән болса,
Тәуелсіздік сорты 1000 дән массасының ең төмен мәнімен ерекшеленді.
Дәннің жеткіліксіз ірілігі өнудің біркелкілігіне және алынатын уыттың
біртектілігіне кері әсер етуі мүмкін.

Зерттелген бидай сорттарының дән шынылығы 55-70\% аралығында өзгерді,
бұл сыра қайнатуға арналған арпаға тән оңтайлы диапазоннан (50-60\%)
жоғары. Әсіресе Матай мен Ертіс сорттарындағы шынылықтың жоғары болуы
эндосперм құрылымының тығыз екенін және ақуызды заттардың мөлшері
жоғарылауын көрсетеді. Бұл уыттандыру кезінде ферментативтік модификация
процестерінің жүруін қиындатып, суландыру мен өндіру режимдерін
қарқындатуды талап етеді.

Зерттелген бидай сорттарының дәніндегі ақуыз мөлшері 11,5-14,0\% құрады.
Ақуыздың ең жоғары мөлшері Ертіс сортында (14,0\%) анықталды, бұл уыттың
экстрактивтілігінің төмендеуіне және суслоның сүзілгіштік қасиеттерінің
нашарлау ықтималдығына байланысты оның уыттық құндылығын төмендетеді.
Ақуыз мөлшері ең төмен (11,5\%) болатын Тәуелсіздік сорты уытқа
қойылатын талаптар тұрғысынан анағұрлым қолайлы болып табылады. Дегенмен
ақуыздың жоғары мөлшері сыраның көбіктүзгіштік қабілетіне оң әсер етуі
мүмкін, бұл сыра өнімдерінің жекелеген түрлерін өндіруде маңызды.

Зерттелген сорттардың барлығы бойынша дәннің физиологиялық күйі
қанағаттанарлық деп бағаланды. Егін жинағаннан кейінгі пісіп-жетілу
кезеңі аяқталған соң өнгіштік қабілеті 95\% болды, ал жаңадан жиналған
дәннің тіршілікке қабілеттілігі кемінде 90\% құрады. Бұл көрсеткіштер
уыттық шикізатқа қойылатын талаптарға сәйкес келеді және біркелкі өнген
уыт алуға мүмкіндік береді.

Зерттеу нәтижелері Матай, Ертіс және Тәуелсіздік бидай сорттарының дәні
уыт өндірісінде қолдануға жеткілікті физиологиялық белсенділікке ие
екенін, алайда технологиялық тұрғыдан маңызды бірқатар көрсеткіштер
бойынша сыра қайнатуға арналған арпадан төмен екенін көрсетеді. Уыт
өсіру мақсатта қолдану үшін ең перспективалы ақуыз мөлшерінің
салыстырмалы түрде төмендігімен және шынылығының орташа деңгейімен
сипатталатын Тәуелсіздік сорты болып табылады. Матай және Ертіс сорттары
эндоспермнің модификациялану үдерістерін қарқындатуға бағытталған уыт
өсіру режимдерін түзетуді қажет етеді. Осыған байланысты сыра қайнатуға
арналған уыт өндіру бойынша кейінгі зерттеулер Тәуелсіздік бидай сорты
негізінде жүргізілді.

Зерттеудің келесі кезеңінде Тәуелсіздік бидай сорты дәнінің
физиологиялық қасиеттерін ескере отырып, уыт өсіру режимдері іріктелді.

Уыт өндіру барысында негізгі бақыланатын параметрлерге уыт өсіру уақыты,
температура және дәндегі ылғалдың массалық үлесі жатады.

Арпамен салыстырғанда, бидайды уыт өсіру үшін суландыру кезеңдерінің
уақыты қысқарақ болып, олардың саны көбірек болуы ұсынылады. Бұл бидай
дәнінде сыртқы қабықшаның болмауына байланысты, соның салдарынан дән
ылғалдың массалық үлесін тез жинақтап та, тез жоғалтып та алады,
сондай-ақ материалдың жабысып-үйілуі мәселелері туындайды. Бидайды уыт
өсіруде жалпы үдерісі арпамен салыстырғанда қысқа уақыт ішінде өтеді
{[}14{]}. Жүргізілген мониторинг нәтижесінде уыт өсірудің белгілі
тәсілдері негізінде 3 режим алынды (1-сурет).

Бірінші уыт өсіру үшін 15ºC температурада 4,5 тәулік бойы жүргізілетін
режим таңдалды. Үдерісті тұрақты температурада жүргізу ұрықтың дамуы мен
дәннің өсуіне әсер ететін күйзелістік факторларды болдырмауға мүмкіндік
береді.

Екінші уыт өсіру жоғары, біртіндеп төмендетілетін температураларда
19ºC-тан бастап, 16ºC-қа дейін біртіндеп азайта отырып, 4 тәулік бойы
жүргізілді. Уыт өсіруді жоғары температураларда жүргізу уыт алудың жалпы
уақытын қысқартуға, өндірістік шығындарды азайтуға мүмкіндік береді.

Үшінші уыт өсіруде төмен, біртіндеп жоғарылайтын температуралар
қолданылды: үдеріс 11ºC-тан басталып, температура 15ºC-қа дейін
біртіндеп көтерілді, өндіру ұзақтығы 5 тәулік болды. Төмен
температураларда уыт өсіру кезінде дәннің ерігіштенуі неғұрлым толық
жүреді, ферменттердің (амилазалардың) жиналуы жақсарады, сондай-ақ
β-глюканның және диметилсульфидтің түзілуі азаяды. Алайда бидайды уыт
өсіруде төмен температурада алынған дайын уыт жоғары температураларда
алынған бидай уытымен салыстырғанда диастатикалық күштің төменірек
деңгейін көрсетеді. Төмен температуралар бидай суслосының тұтқырлығын
төмендетуге мүмкіндік береді.

Бидайды өндіру кезінде күн сайын келесі көрсеткіштер бақылауға алынды:
ылғалдылық (2-сурет), дәндегі азотты заттар деңгейі (3-сурет), бос
аминді азот деңгейі (4-сурет).

2-суретте өндіру барысында ылғал мөлшерінің өзгеруі көрсетілген.
\end{multicols}

\fig[0.6\textwidth]{p2/image24}[1 - сурет. Уыт өсурідің температуралық
режимдері]

\fig[0.6\textwidth]{p2/image25}[2 - сурет. Уыт өсіру барысында
ылғалдылық мөлшерінің өзгеруі]

\begin{multicols}{2}
2-суреттен екінші суландырудан кейін (24 сағаттан соң) бидайдың ылғалды
тез жоғалта бастайтыны көрінеді. Бұл үдеріс 1 сағат ішінде шамамен
\textasciitilde1\% ылғал жинау жылдамдығымен бірдей қарқынмен жүреді.
Сондықтан жиі су бүрку жүргізілді - тәулігіне 4--5 рет.

Бидай - қауыссыз дақыл, сондықтан ол арпаға қарағанда ылғалды тезірек
сіңіреді және тезірек қайта жоғалтады; осыған байланысты бидай үшін
үшінші жібіту кезеңі қолданылмайды.2-суреттен дәнді суға батырған кезде
ылғал мөлшерінің күрт артатыны байқалады.1-режим және 2-режим кезінде
ылғалдың 36\% және 41\%-ға дейін (тиісінше) күрт өсуі уыт өсірудің
18-сағатында тіркелді.3-режимде (уыт өсіру температурасы төмен
болғанда) 18-сағатта ылғал деңгейінің күрт көтерілуі байқалмады: дәннің
ылғалдылығы бірқалыпты артып, ісінуі және өнуі жүрді.

Үш түрлі режимнің графиктерін салыстырғанда, дәндегі ылғалдың 40\%
деңгейіне жетуі уақыт бойынша ығысумен өтті: 2-режимде (жоғары
температурада) - 28 сағатта, 1-режимде (стандартты температурада) -
шамамен 40 сағатта, 3-режимде (төмен температурада) - 48 сағатта.
Уыттағы ылғалдың ең жоғары мөлшері 3-режимді (төмен температураларда)
қолданғанда алынды.

Кезең-кезеңімен су бүрку кезінде дәндегі ылғал мөлшері қатты құбылады:
су бүркуден кейін ылғалдылық тез өсіп, өндіру барысында біртіндеп
төмендейді. Дәнді суға батыруға негізделген тәсілдер бидайды уыт өсіру
үшін қолайсыз. Демек, бидайды өндіру үшін үздіксіз су бүрку арқылы
өндіру әдісі анағұрлым тиімді болып табылады.

3-суретте өндіру барысында дәндегі азотты заттар мөлшерінің өзгерісі
көрсетілген.
\end{multicols}

\fig[0.6\textwidth]{p2/image26}[3 - сурет. Уыт өсіру барысында дәндегі
азотты заттар мөлшерінің өзгеруі]
\fig[0.6\textwidth]{p2/image27}[4 - сурет. Уыт өсіру үдерісі кезінде
аминді азоттың мөлшерінің өзгеруі]

\begin{multicols}{2}
3-диаграммадан дәндегі азотты заттардың үлесі үнемі құбылып отыратыны
көрінеді. Азотты заттардың төмендеуі дәндегі тыныс алу үдерістері,
тамыршалар мен өскіндердің өсуі, сондай-ақ суландыру кезінде ішінара
шайылып кетуі есебінен жүрді. Сыра қайнатуға арналған арпада азотты
заттар мөлшерінің төмендеуі уыт өсірудің шамамен 4-тәулігінде байқалса,
бидайда бұл төмендеу шамамен 2-тәулікте көрінеді. Осы кезеңде суландыру
үдерісі аяқталды, яғни үдерістің басынан аяқталғанға дейін азотты
заттардың шайылуы орын алды. Төмен температурадағы үшінші режимде
суландыру ең ұзақ болды, соған байланысты азотты заттардың азаюы да ең
жоғары деңгейде байқалды, бұл 3-суретте расталады.

Бос аминді азот деңгейі бойынша ақуыздардың гидролизін бағалауға болады:
бұл көрсеткіш қаншалықты жоғары болса, соншалықты көп ақуыз ыдырағанын
білдіреді. Бос аминді азот ашытқылар үшін негізгі қоректік көз болып
табылады. Бидай уыты үшін арпа уытымен салыстырғанда ақуыздардың азырақ
қарқынды ыдырауы қолайлы, себебі бұл бидай уытына тән хош иісті
қалыптастыруға мүмкіндік береді. Бидай уытында бос аминді азоттың
мөлшері шамамен 90--120 мг/100 г құрғақ затқа (ҚЗ) тең болуы тиіс.
4-суретте уыт өсіну барысында аминді азот мөлшерінің өзгеруі
көрсетілген.

Өндіру барысында аминді азот мөлшерінің өзгерісін зерттеу уыт өсіру
температурасына тікелей тәуелділікті көрсетті (4-сурет). Уыт өсіру
үдерісін жоғары температурада жүргізгенде және соған сәйкес қысқа
мерзімде (2-режим) аминді азоттың жиналуы төмен температураларда әрі
ұзақ уыттандыру кезінде (3-режим) қарағанда жылдамырақ жүрді.

Уыт өсірудің үш түрлі режимін зерттеу нәтижелері ең оңтайлысы 2-режим
(жоғары, біртіндеп төмендетілетін температуралардағы уыттандыру) екенін
көрсетті. Осы режим бойынша алынған уыт үлгісінде ақуыз мөлшері ең төмен
болды және амилолитикалық ферменттердің ең көп мөлшері түзілді
(көрсеткіші: диастатикалық күш). Бұл уыт өсіру технологиясы дайын уыттың
сапасын жақсартатын параметрлерді ғана емес, сонымен қатар уыт өсірудің
қысқарақ уақытын да қамтиды.

{\bfseries Қорытынды.} Осы зерттеудің мақсаты Қазақстан Республикасының
отандық селекциялы бидай сорттарынан сыра қайнату өнеркәсібінің
қажеттіліктері үшін бидай уытын өндірудің мақсатқа сәйкестігін негіздеу
болды. Қойылған мақсатқа жету үшін ғылыми-техникалық және нормативтік
әдебиеттерді талдау мен жалпылау әдістері, сондай-ақ дән мен уыт сапасын
бағалаудың жалпы қабылданған физика-химиялық, биохимиялық және
технологиялық әдістері қолданылды. Эксперименттік бөлім Тәуелсіздік
бидай сортының әртүрлі температуралық-уақыттық уыт өсіру режимдерін
іріктеу және салыстыруды, сондай-ақ ылғал жинақталуы динамикасын, масса
жоғалтуларын, азотты заттар мөлшерін және бос аминді азот деңгейін
бақылауды қамтыды.

Жүргізілген зерттеулер нәтижесінде Матай, Ертіс және Тәуелсіздік бидай
сорттарының дәні жоғары физиологиялық белсенділікпен және
қанағаттанарлық өнгіштік көрсеткіштерімен сипатталатыны анықталды, бұл
оны уыт өндіруге арналған әлеуетті шикізат ретінде қарастыруға мүмкіндік
береді. Сонымен қатар, сорттардың уыттық құндылығына әсер ететін
физика-химиялық сипаттамалары бойынша елеулі айырмашылықтар белгіленді.
Ең перспективалысы ақуыз мөлшері мен шынылығы орташа деңгейде болатын
Тәуелсіздік сорты болып шықты, бұл эндоспермнің ферментативтік
модификациялануына қолайлырақ жағдай жасайды. Уыт өсіру режимдерін
эксперименттік іріктеу температура мен өндіру ұзақтығы ылғал
жинақталуының динамикасына және ақуызды заттардың гидролизіне шешуші
әсер ететінін көрсетті. Жоғары температуралық режимдер қысқа уақыт
ішінде бос аминді азоттың қарқынды жиналуына ықпал етсе, төмен
температуралық уыт өсіру дәннің анағұрлым біркелкі модификациялануын
және бидай суслосының тұтқырлығының төмендеуін қамтамасыз етті.
Сондай-ақ, бидай үшін арпаға тән көп мәрте суландырудың болдырмайтын,
үздіксіз су бүрку арқылы өндіру әдісі қолайлы екені анықталды.

Алынған деректер негізінде қазақстандық бидай сорттарынан сыра қайнату
өнеркәсібінің талаптарына сәйкес келетін ферментативтік белсенділік және
азотты алмасу көрсеткіштері бар бидай уытын алудың принциптік мүмкіндігі
туралы қорытындылар жасалды. Уыт өсіру режимдерін оңтайландыру
жағдайында отандық селекциялы бидай дәні алдын ала белгіленген
технологиялық қасиеттері бар уыттың қалыптасуын қамтамасыз ете алады
деген ұсынылған гипотеза расталды. Сонымен бірге шынылығы және ақуыз
мөлшері жоғары сорттардың модификация үдерістерін қарқындатуға
бағытталған технологиялық параметрлерді түзетуді қажет ететіні
көрсетілді.

Зерттеу нәтижелерін тәжірибелік тұрғыдан енгізудің мүмкіндіктері
Қазақстанның сыра қайнату саласының шикізат базасын жергілікті дәнді
шикізатты пайдалану есебінен кеңейтумен, импорттық бидай уытына
тәуелділікті төмендетумен және өнімнің өзіндік құнын оңтайландырумен
байланысты.

Дайындалған уыт өсіру режимдері бидай уытын өнеркәсіптік көлемде
өндіруде қолданылуы мүмкін, сондай-ақ уыт өсіру технологияларын
жетілдіруге, уыттың оңтайлы түрлерін іріктеуге және алдын ала
белгіленген органолептикалық сипаттамалары бар бидай сырасының жаңа
түрлерін жасауға бағытталған кейінгі зерттеулерге негіз бола алады.

\emph{Мақаланы дайындау барысында жасанды интеллект құралдары, соның
ішінде OpenAI ұсынған GPT-4 тілдік моделі қолданылды. Модель тек мәтінді
редакциялауға және ғылыми экспозицияны құрылымдауға көмектесу үшін
қолданылды. Барлық нәтижелер автордың жеке еңбегімен алынған.}
\end{multicols}

\begin{center}
{\bfseries Әдебиеттер}
\end{center}

\begin{refs}
1. Кунце В., Мит Г. Технология солода и пива: пер. с нем. - СПб.:
Профессия, 2001. - 912 с.- ISBN 5-93913-006-2.

2. Briggs D.E., Boulton C.A., Brookes P.A., Stevens R. Brewing: Science
and Practice. - Cambridge: Woodhead Publishing; Boca Raton: CRC Press,
2004. - 882 p. ISBN 978-1-85573-490-6.

3. Kunze W. Technology Brewing \& Malting. - 5th revised English ed. -
Berlin: VLB Berlin, 2014. - ISBN 978-3-921690-77-2.

4. Нарцисс Л. Пивоварение. Т.1: Технология солодоращения: пер. с нем. /
под общ. ред. Г. А. Ермолаевой, Е. Ф. Шаненко; пер. А. С. Яблоковой. -
СПб.: Профессия, 2007. - 584 с. ISBN 5-93913-118-2.

5. Меледина Т.В., Прохорчик И.П., Кузнецова Л.И. Биохимические процессы
при производстве солода - СПб.: НИУ ИТМО; ИХиБТ, 2013. - 89 с.

6. Stewart G.G., Priest F.G. (eds.). Handbook of Brewing. - 2nd ed.-
Boca Raton: CRC Press, 2006. - 872 p. ISBN 0-8247-2657-X

7. Bamforth C.W. Beer: Tap into the Art and Science of Brewing. - 3rd
ed.-Oxford: Oxford University Press, 2009. - ISBN 978-0-19-530542-5.

8. FAO. FAOSTAT: Production: Crops and livestock products (Wheat):
dataset. - Rome: Food and Agriculture Organization of the United
Nations, {[}online{]}. -URL: data.un.org/Data.aspx?
d=FAO\&f=itemCode\%3A15 - (accessed on 15.12.2025).

9. Bamforth C. W. (ed.). Beer: A Quality Perspective. - Amsterdam;
Boston: Elsevier/Academic Press, 2009. - ISBN 978-0-12-669201-3.

10. Faltermaier A., Waters D., Becker T., Arendt E.K., Gastl M. Common
wheat (Triticum aestivum L.) and its use as a brewing cereal - a review
// Journal of the Institute of Brewing. - 2014.- Vol.120(1). - P.1-15.
DOI 10.1002/jib.107.

11. Poreda A., Bijak M., Zdaniewicz M., Jakubowski M., Makarewicz M.
Effect of wheat malt on the concentration of metal ions in wort and
brewhouse by-products // Journal of the Institute of Brewing. - 2015. -
Vol.121(2). - P.224-230. DOI 10.1002/jib.226.

12. ASBC. Methods of Analysis: online edition. - American Society of
Brewing Chemists, {[}online{]}. - URL: asbcnet.org/methods. - (accessed
on 15.12.2025)

13. Меледина Т.В. Сырье и вспомогательные материалы в пивоварении. -
СПб.: Профессия, 2003. - 350 с. ISBN 5-93913-054-2.

14. Bakytzhan D., Askarbekov E., Klupsaite D., Mockus E., Starkute V.,
Mozuriene E., Ruibys R., Bartkiene E. Effects of different wheat malt on
beer characteristics, including amino acids, gamma aminobutyric acid,
biogenic amines, and volatile compounds // Applied Food Research. -
2025. -Vol.5(2):101203. DOI 10.1016/j.afres.2025.101203.
\end{refs}

\begin{center}
{\bfseries References}
\end{center}

\begin{refs}
1. Kunce V., Mit G. Tehnologija soloda i piva: per. s nem. - SPb.:
Professija, 2001.- 912 s.- ISBN 5-93913-006-2. {[}in Russian{]}

2. Briggs D.E., Boulton C.A., Brookes P.A., Stevens R. Brewing: Science
and Practice. - Cambridge: Woodhead Publishing; Boca Raton: CRC Press,
2004. - 882 p. ISBN 978-1-85573-490-6.

3. Kunze W. Technology Brewing \& Malting. - 5th revised English ed. -
Berlin: VLB Berlin, 2014. - ISBN 978-3-921690-77-2.

4. Narciss L. Pivovarenie. T.1: Tehnologija solodorashhenija: per. s
nem. / pod obshh. red. G. A. Ermolaevoj, E. F. Shanenko; per. A. S.
Jablokovoj. - SPb.: Professija, 2007. - 584 s. ISBN 5-93913-118-2. {[}in
Russian{]}

5. Meledina T.V., Prohorchik I.P., Kuznecova L.I. Biohimicheskie
processy pri proizvodstve soloda - SPb.: NIU ITMO; IHiBT, 2013. - 89 s.
{[}in Russian{]}

6. Stewart G.G., Priest F.G. (eds.). Handbook of Brewing. - 2nd ed.-
Boca Raton: CRC Press, 2006. - 872 p. ISBN 0-8247-2657-X

7. Bamforth C.W. Beer: Tap into the Art and Science of Brewing. - 3rd
ed.-Oxford: Oxford University Press, 2009. - ISBN 978-0-19-530542-5.

8. FAO. FAOSTAT: Production: Crops and livestock products (Wheat):
dataset. - Rome: Food and Agriculture Organization of the United
Nations, {[}online{]}. -URL: data.un.org/Data.aspx?
d=FAO\&f=itemCode\%3A15 - (accessed on 15.12.2025).

9. Bamforth C. W. (ed.). Beer: A Quality Perspective. - Amsterdam;
Boston: Elsevier/Academic Press, 2009. - ISBN 978-0-12-669201-3.

10. Faltermaier A., Waters D., Becker T., Arendt E.K., Gastl M. Common
wheat (Triticum aestivum L.) and its use as a brewing cereal - a review
// Journal of the Institute of Brewing. - 2014.- Vol.120(1). - P.1-15.
DOI 10.1002/jib.107.

11. Poreda A., Bijak M., Zdaniewicz M., Jakubowski M., Makarewicz M.
Effect of wheat malt on the concentration of metal ions in wort and
brewhouse by-products // Journal of the Institute of Brewing. - 2015. -
Vol.121(2). - P.224-230. DOI 10.1002/jib.226.

12. ASBC. Methods of Analysis: online edition. - American Society of
Brewing Chemists, {[}online{]}. - URL: asbcnet.org/methods. - (accessed
on 15.12.2025)

13. Meledina T.V. Syr' e i
vspomogatel' nye materialy v pivovarenii. - SPb.:
Professija, 2003.- 350 s. ISBN 5-93913-054-2. {[}in Russian{]}

14. Bakytzhan D., Askarbekov E., Klupsaite D., Mockus E., Starkute V.,
Mozuriene E., Ruibys R., Bartkiene E. Effects of different wheat malt on
beer characteristics, including amino acids, gamma aminobutyric acid,
biogenic amines, and volatile compounds // Applied Food Research. -
2025. -Vol.5(2):101203. DOI 10.1016/j.afres.2025.101203.
\end{refs}

\begin{info}
\emph{{\bfseries Авторлар туралы мәліметтер}}

Бакытжан Д. Н. - Алматы технологиялық университетінің, Алматы,
Қазақстан, е-mail: didar\_99\_kz@mail.ru;

Аскарбеков Э.Б. - PhD, қауымдастырылған профессор. Қ.Құлажанов
атындағы Қазақ технология және бизнес университеті, Астана, Қазақстан,
е-mail: erik\_ab82@mail.ru;

Новикова И. В. - т. ғ. д., доцент, ФГБОУ ВО "Воронежский
государственный университет инженерных технологий", Ресей Федерациясы,
е-mail: noviv@list.ru;

Кекибаева А.К. - PhD, Алматы технологиялық университеті, Алматы,
Қазақстан, е-mail: kekibaevaa@gmail.com;

Байгазиева Г. И. - б. ғ. к., профессор, Алматы технологиялық
университеті, Алматы, Қазақстан, е-mail: gulgaisha.atu@gmail.com.

\emph{{\bfseries Information about the authors}}

Bakytzhan D. N. - Almaty Technological University, Almaty, Kazakhstan,
e-mail: didar\_99\_kz@mail.ru:

Askarbekov E. B. - PhD, associate professor. K. Kulazhanov Kazakh
University of technology and business,Astana, Kazakhstan, e-mail:
erik\_ab82@mail.ru:

Novikova I. V. - Doctor of technical sciences, associate professor,
FSBEI HE "Voronezh State University of engineering technologies",
Russian Federation, e-mail: noviv@list.ru;

Kekibayeva A. K. - PhD, Almaty Technological University, Almaty,
Kazakhstan, e-mail: kekibaevaa@gmail.com;

Baigazieva G. I - candidate of biological sciences, professor, Almaty
Technological University, Almaty, Kazakhstan, e-mail:
gulgaisha.atu@gmail.com.
\end{info}
